\documentclass[12pt,a4paper]{article}

\usepackage[utf8]{inputenc}
\usepackage[T1]{fontenc}
\usepackage[brazil]{babel}
\usepackage{lmodern}
\usepackage{geometry}
\usepackage{microtype}
\usepackage{booktabs}
\usepackage{longtable}
\usepackage[hidelinks]{hyperref}
\usepackage{listings}
\usepackage{xcolor}

\geometry{margin=2.5cm}
\setlength{\parskip}{0.45em}
\setlength{\parindent}{1.25cm}
\selectlanguage{brazil}

\lstdefinestyle{terminal}{
  basicstyle=\ttfamily\small,
  breaklines=true,
  frame=single,
  columns=fullflexible,
  keepspaces=true,
  backgroundcolor=\color{gray!6}
}

\title{\textbf{Relatório técnico do Video Review}\\
\large Pipeline experimental para rastreamento e detecção de instrumentos cirúrgicos}
\author{Projeto acadêmico -- Tech Challenge Fase 4}
\date{\today}

\begin{document}
\maketitle

\begin{abstract}
Este relatório descreve a implementação \textit{video-review}, aplicação experimental para selecionar e rastrear instrumentos em vídeos de colecistectomia, preparar anotações no formato YOLO, treinar um detector de objetos e executar inferência em vídeo. A implementação usa Python, OpenCV e Ultralytics YOLOv8, organizando os casos de uso em uma camada de aplicação separada dos adaptadores de infraestrutura. A preparação atual divide aleatoriamente as imagens por ferramenta nas proporções 80\% para treino, 10\% para validação e 10\% para teste. Como imagens próximas do mesmo vídeo podem cair em partições diferentes, os resultados não demonstram generalização para procedimentos independentes. Os artefatos de treino fornecidos para este relatório indicam mAP@50 de 0,974 e mAP@50--95 de 0,894 na validação; esses valores são experimentais e não representam avaliação no conjunto de teste ou validação clínica. O protótipo destina-se exclusivamente à revisão humana e não diagnostica nem orienta decisões cirúrgicas.
\end{abstract}

\noindent\textbf{Palavras-chave:} visão computacional; rastreamento de objetos; YOLOv8; OpenCV; instrumentos cirúrgicos; revisão humana.

\section{Introdução}

Vídeos cirúrgicos contêm instrumentos parcialmente ocluídos, em movimento e sujeitos a mudanças de escala e aparência. Este projeto implementa um fluxo experimental para gerar anotações espaciais a partir de rastreamento visual, treinar um detector de ferramentas e produzir resultados que possam ser inspecionados por uma pessoa. O escopo é compatível com a proposta registrada em \texttt{documentos/requisitos/prototipo-revisao-cirurgica.md}: apoio acadêmico à revisão posterior, sem alegações clínicas.

O sistema não considera uma detecção ausente como evidência de ausência do objeto e não trata automaticamente objetos desconhecidos como anomalias. Deriva do rastreador, oclusão, perda da ferramenta e qualidade das caixas são limitações que exigem inspeção humana.

\section{Objetivos e escopo}

O objetivo técnico é implementar um pipeline reprodutível capaz de:
\begin{enumerate}
  \item selecionar manualmente uma região inicial de uma ferramenta em um vídeo;
  \item propagar essa região por meio de um tracker do OpenCV;
  \item salvar frames e labels espaciais compatíveis com YOLO;
  \item validar e preparar um dataset dividido por ferramenta;
  \item treinar YOLOv8 com acompanhamento de métricas na validação; e
  \item executar inferência em vídeo, salvando visualização e dados tabulares.
\end{enumerate}

A implementação contempla sete classes: \texttt{grasper}, \texttt{bipolar}, \texttt{hook}, \texttt{scissors}, \texttt{clipper}, \texttt{irrigator} e \texttt{specimen\_bag}. A ordem define os IDs YOLO de 0 a 6.

\section{Arquitetura e organização}

O pacote principal está em \texttt{src/video\_review/} e separa a interface de linha de comando, os casos de uso e os adaptadores externos:

\begin{lstlisting}[style=terminal]
src/video_review/
|-- cli.py
|-- models.py
|-- application/
|   |-- track_video.py
|   |-- dataset.py
|   |-- train_detector.py
|   `-- predict_video.py
`-- infrastructure/
    |-- opencv_video_tracker.py
    |-- ultralytics_trainer.py
    `-- ultralytics_video_predictor.py
\end{lstlisting}

\subsection{Camada de interface e modelos}

O módulo \texttt{cli.py} define os comandos \texttt{track}, \texttt{data validate}, \texttt{data prepare}, \texttt{detector train} e \texttt{detector predict}. \texttt{models.py} contém os objetos de solicitação e resultado do rastreamento. Os casos de uso de treino e inferência também declaram objetos próprios para seus parâmetros e saídas.

\subsection{Camada de aplicação}

\texttt{application/track\_video.py} verifica a existência do vídeo, o tracker e a classe antes de delegar a execução. \texttt{application/dataset.py} lê os rastreamentos por pasta de ferramenta, valida formato e coordenadas dos labels, realiza a divisão e escreve o dataset YOLO. \texttt{application/train\_detector.py} valida os argumentos básicos e coordena o treinador. \texttt{application/predict\_video.py} valida vídeo, pesos e limiar de confiança antes de delegar a inferência.

\subsection{Infraestrutura}

\texttt{opencv\_video\_tracker.py} implementa a leitura do vídeo, seleção interativa do frame e da caixa inicial, rastreamento CSRT/KCF e gravação das imagens e labels. \texttt{ultralytics\_trainer.py} chama o treino do Ultralytics e retorna o caminho de \texttt{best.pt}. \texttt{ultralytics\_video\_predictor.py} executa predições em fluxo, grava um MP4 anotado e exporta detecções para CSV.

Essa organização isola os fluxos da aplicação dos detalhes específicos de OpenCV e Ultralytics, sem exigir que os casos de uso conheçam diretamente suas APIs.

\subsection{Pastas de dados e artefatos}

\begin{longtable}{p{0.36\textwidth}p{0.56\textwidth}}
\toprule
\textbf{Caminho} & \textbf{Uso} \\
\midrule
\endfirsthead
\toprule
\textbf{Caminho} & \textbf{Uso} \\
\midrule
\endhead
\texttt{resources/cholec80/videos/} & Vídeos de entrada disponibilizados localmente. \\
\texttt{resources/cholec80/tracking/} & Imagens e labels de rastreamentos, organizados por ferramenta e execução. \\
\texttt{resources/datasets/cholec80-yolo/} & Dataset YOLO preparado, manifesto e configuração; destino padrão recriado pelo preparador. \\
\texttt{runs/detect/} & Pesos, métricas, gráficos e configuração de cada treino. \\
\texttt{resources/results/} & Saídas padrão da inferência em vídeo. \\
\texttt{tests/} & Testes automatizados, incluindo fixtures temporárias para preparação. \\
\bottomrule
\end{longtable}

\section{Método e fluxo de processamento}

\subsection{Rastreamento e criação de anotações}

O operador seleciona uma ferramenta, navega até um frame apropriado e desenha uma caixa delimitadora inicial. O tracker CSRT ou KCF propaga a região nos frames subsequentes. Para cada frame, o sistema salva a imagem original de treino e um arquivo de label YOLO com ID da classe, centro normalizado e largura/altura normalizadas. Quando o rastreamento falha, o código salva label vazio; esses frames podem funcionar como exemplos sem caixa e também devem ser revisados antes do uso.

Cada execução é armazenada sob a pasta da classe e recebe um identificador sequencial associado ao vídeo. Os rastreamentos de ferramentas diferentes são independentes. Assim, coincidência de número de vídeo e índice de frame em duas pastas não autoriza fundir as imagens nem combinar seus labels.

\subsection{Validação e preparação do dataset}

O validador exige imagens JPEG com labels correspondentes, rejeita labels órfãos, verifica os cinco campos do formato YOLO, confere o ID com a pasta da ferramenta e verifica que coordenadas sejam numéricas, finitas e pertençam ao intervalo normalizado \([0,1]\).

A preparação atual realiza sorteio aleatório reprodutível separadamente para cada ferramenta, com seed configurável e proporções padrão 80\%/10\%/10\%. A unidade amostral é o registro do rastreamento, não a chave global vídeo-frame; portanto, imagens de rastreamentos independentes não são colapsadas. Arredondamentos inteiros podem causar pequenas diferenças em relação às proporções exatas, sobretudo em classes com poucas imagens. O preparador exige amostras suficientes para preencher as três partições.

O destino é removido e recriado antes da cópia; esse comportamento limita-se ao caminho de saída configurado e protegido por verificações contra destinos inseguros. A saída inclui \texttt{data.yaml}, pastas \texttt{images/} e \texttt{labels/} para treino, validação e teste, e \texttt{split\_manifest.json}. O manifesto registra estratégia, seed, proporções, contagens por partição e ferramenta e a lista de arquivos alocados.

\paragraph{Limitação temporal.} O sorteio por imagem não considera a distância entre índices de frame. Consequentemente, frames adjacentes ou visualmente semelhantes podem aparecer em partições distintas. As métricas podem, portanto, superestimar o desempenho em novos dados do mesmo vídeo. A estratégia atual atende à divisão estratificada acordada para o experimento, mas não equivale a separação por procedimento e não mede generalização para vídeos totalmente independentes.

\subsection{Treinamento}

O treinamento usa o checkpoint inicial indicado (por padrão, \texttt{yolov8n.pt}), os frames de treino e a validação definida no \texttt{data.yaml}. O Ultralytics seleciona e salva o melhor checkpoint em \texttt{weights/best.pt}, além do último em \texttt{weights/last.pt}, e produz \texttt{results.csv}, gráficos e argumentos da execução. A pasta padrão é \texttt{runs/detect/cholec80-yolov8n/}; execuções com nome existente são direcionadas a uma nova pasta pelo Ultralytics, pois \texttt{exist\_ok=False}.

Os parâmetros de tamanho de imagem e batch influenciam recursos e custo. \texttt{imgsz=640} redimensiona imagens para o tamanho de entrada do modelo; aumentar pode preservar detalhes de ferramentas pequenas, com maior custo de memória e processamento. \texttt{batch=2} processa duas imagens por batch antes da atualização; batch maior pode melhorar utilização do dispositivo, mas consome mais memória. A configuração usada deve ser registrada ao comparar experimentos.

O teste não deve participar do ajuste do modelo nem da escolha de hiperparâmetros. A avaliação final no split \texttt{test} deve ser realizada somente após congelar as decisões tomadas com treino e validação.

\subsection{Inferência}

Na inferência, o modelo processa o vídeo em fluxo e aplica o limiar de confiança configurado (padrão 0,25). Para cada frame, salva as caixas desenhadas no vídeo de saída. O CSV registra índice do frame, timestamp calculado a partir do FPS, ID e nome da classe, confiança e coordenadas \((x_1,y_1,x_2,y_2)\). Se os arquivos de saída já existirem, a aplicação recusa sobrescrevê-los.

\section{Execução}

As dependências de aplicação estão listadas em \texttt{requirements.txt}. Exemplo de instalação e comandos principais:

\begin{lstlisting}[style=terminal]
python3 -m venv .venv
source .venv/bin/activate
python -m pip install -r requirements.txt

PYTHONPATH=src python3 -m video_review.cli data validate
PYTHONPATH=src python3 -m video_review.cli data prepare --seed 42

PYTHONPATH=src python3 -m video_review.cli detector train \
  --dataset resources/datasets/cholec80-yolo/data.yaml \
  --model yolov8n.pt --epochs 50 --imgsz 640 --batch 2 \
  --name cholec80-yolov8n
\end{lstlisting}

Para inspecionar as opções disponíveis, usar \texttt{PYTHONPATH=src python3 -m video\_review.cli detector train --help}. Os comandos de rastreamento e inferência requerem interface gráfica e/ou arquivos de vídeo e pesos acessíveis localmente.

\section{Resultados experimentais disponíveis}

Foi compartilhado um treino de 50 épocas, com modelo YOLOv8 nano, tamanho de imagem 640, batch 2 e validação ativada. O diretório indicado em \texttt{args.yaml} é \texttt{runs/detect/cholec80-yolov8n-2}. O CSV de resultados apresenta na época 50 precisão 0,894, recall 0,999, mAP@50 0,966 e mAP@50--95 0,889. O maior mAP@50--95 no CSV é 0,894 na época 47. A curva de precisão-recall reportada no gráfico indica mAP@50 agregado de 0,974 e AP de 0,863 para \texttt{irrigator}, inferior às demais classes mostradas.

Esses números são da validação, não do teste. O próprio gráfico da matriz de confusão mostra falsos positivos na categoria de fundo, com maior volume para algumas classes; a classe \texttt{irrigator} também apresenta uma curva de confiança--F1 menos favorável. A inspeção de imagens, caixas e falsos positivos continua necessária. Não foram fornecidos resultados de avaliação final no conjunto de teste para este relatório.

\begin{table}[ht]
\centering
\caption{Métricas agregadas da última época no arquivo de resultados compartilhado.}
\begin{tabular}{lr}
\toprule
\textbf{Métrica de validação} & \textbf{Valor na época 50} \\
\midrule
Precisão & 0,894 \\
Recall & 0,999 \\
mAP@50 & 0,966 \\
mAP@50--95 & 0,889 \\
\bottomrule
\end{tabular}
\end{table}

O valor mAP@50--95 da melhor época no CSV foi 0,894 na época 47; a seleção do Ultralytics para \texttt{best.pt} pode considerar sua métrica de fitness composta. Portanto, os números da época 50 não devem ser confundidos automaticamente com a avaliação dos pesos \texttt{best.pt}.

\section{Testes e reprodutibilidade}

Os testes do pipeline são executados com:

\begin{lstlisting}[style=terminal]
PYTHONPATH=src python3 -m unittest discover -s tests -v
\end{lstlisting}

Os testes de preparação verificam a estratificação, a separação de registros com mesmo vídeo/frame em rastreamentos diferentes, a limpeza do destino configurado, validação de labels e comportamento determinístico. Eles usam diretórios temporários e não recriam o dataset real. A seed de preparação deve ser registrada junto aos resultados para permitir reconstrução da divisão, desde que a fonte não mude.

\section{Limitações, segurança e validade}

\begin{itemize}
  \item A qualidade do detector depende da revisão das caixas rastreadas. Um tracker visual pode derivar, perder o objeto ou incluir fundo.
  \item Labels vazios representam frames sem caixa de objeto no arquivo, inclusive frames após falha de rastreamento; não garantem ausência visual do instrumento.
  \item O sorteio por imagem pode misturar frames temporalmente próximos entre treino, validação e teste, comprometendo a independência estatística das partições.
  \item A distribuição de classes e a baixa quantidade de dados de algumas ferramentas tornam métricas por classe mais incertas.
  \item Vídeo anotado e CSV são artefatos de inspeção; confiança do modelo não é probabilidade clínica calibrada.
  \item Não há, na implementação descrita, módulo validado de detecção de anomalias ou classificação de objetos fora das sete classes.
  \item O protótipo é acadêmico e de apoio à revisão humana. Não é dispositivo médico, não diagnostica e não deve ser usado para decisão intraoperatória.
\end{itemize}

\section{Conclusão e próximos passos}

O \textit{video-review} integra rastreamento interativo, validação/preparação de caixas, treinamento YOLOv8 e inferência em vídeo numa aplicação modular. O experimento compartilhado demonstra métricas de validação promissoras para o split utilizado, mas a divisão aleatória por imagem e a possibilidade de frames vizinhos em conjuntos distintos impedem interpretar esses resultados como evidência de generalização independente.

Como próximos passos técnicos, recomenda-se revisar amostras e falsos positivos por classe, confirmar a composição do manifesto do dataset usado no treino, avaliar uma única vez no conjunto de teste reservado e, para uma estimativa mais rigorosa de generalização, realizar uma avaliação adicional separada por vídeo/procedimento. Qualquer alteração de protocolo deve ser documentada e aplicada sem reutilizar o teste como conjunto de ajuste.

\section*{Referências}
\addcontentsline{toc}{section}{Referências}

\begin{enumerate}
  \item TWINANDA, A. P. et al. EndoNet: A Deep Architecture for Recognition Tasks on Laparoscopic Videos. \textit{IEEE Transactions on Medical Imaging}, 2017.
  \item ULTRALYTICS. \textit{Ultralytics YOLO Documentation}. Disponível em: \url{https://docs.ultralytics.com/}. Acesso em: \today.
  \item OPENCV. \textit{OpenCV Documentation}. Disponível em: \url{https://docs.opencv.org/}. Acesso em: \today.
  \item PROJETO VIDEO-REVIEW. \textit{Protótipo acadêmico de apoio à revisão cirúrgica}. Requisitos locais, \texttt{documentos/requisitos/prototipo-revisao-cirurgica.md}.
\end{enumerate}

\end{document}
